\startcomponent appendix
\product fduthesis

\title{附录 A\quad 补充实验数据}

以下为正文中未完全展示的补充实验数据，包括不同超参数组合下模型的训练结果。

\placetable
  [here][tab:appendix-a]
  {不同超参数组合下的实验结果}
  {\starttabulate[|cw(2.8cm)|cw(2.5cm)|cw(2.5cm)|cw(3cm)|]
    \HL
    \NC {\bf 学习率} \NC {\bf 批大小} \NC {\bf 迭代轮数} \NC {\bf 准确率} \NC \NR
    \HL
    \NC $1\times10^{-3}$  \NC 32  \NC 50  \NC 0.910 \NC \NR
    \NC $1\times10^{-3}$  \NC 64  \NC 50  \NC 0.923 \NC \NR
    \NC $5\times10^{-4}$  \NC 32  \NC 80  \NC 0.935 \NC \NR
    \NC $5\times10^{-4}$  \NC 64  \NC 80  \NC 0.941 \NC \NR
    \NC $1\times10^{-4}$  \NC 32  \NC 100 \NC 0.948 \NC \NR
    \NC $1\times10^{-4}$  \NC 64  \NC 100 \NC 0.962 \NC \NR
    \HL
  \stoptabulate}

由实验结果可以看出，较小的学习率配合较多的迭代轮数能够获得更好的模型性能。
当学习率为 $1\times10^{-4}$、批大小为 64、迭代轮数为 100 时，
模型达到最优准确率 0.962。

\blank[2*line]

\title{附录 B\quad 算法伪代码}

以下为本文所提出的模型训练流程的伪代码描述。

\starttyping
Algorithm: Model Training
Input:  Dataset D, learning rate lr,
        iterations T
Output: Trained parameters theta

1. Initialize theta randomly
2. for t = 1 to T do
3.   Sample mini-batch B from D
4.   Compute loss L(theta; B)
5.   Compute gradient g = grad L(theta)
6.   Update theta = theta - lr * g
7. end for
8. return theta
\stoptyping

\blank[line]

以下为模型推理（预测）流程的伪代码。

\starttyping
Algorithm: Model Inference
Input:  Test sample x,
        trained parameters theta
Output: Predicted label y

1. Compute features h = f(x; theta)
2. Compute logits z = W * h + b
3. Compute probabilities p = softmax(z)
4. return y = argmax(p)
\stoptyping

\stopcomponent
